\section{Renewal targets and reliable realization}\label{sec:renewal}
\subsection{Encoding a real in gap probabilities}
Fix $\beta\in\D_3\cap[0,1]$ and uniformly $0'$-computable
rationals $b_k\in[0,1]$ converging to it. Let the number $G$ of
zeros between successive ones have distribution
\begin{equation}\label{eq:gaplaw}
 p_{2k}=2^{-k}b_k,\qquad
 p_{2k+1}=2^{-k}(1-b_k),\qquad k\geq1.
\end{equation}
The pair weights sum to one. There is a stationary renewal measure
$\nu_b$ with these independent gaps: choose the block containing
the origin by the length-biased distribution of $G+1$, place the
origin uniformly inside that block, and extend by independent gaps
on both sides. This construction is valid even for a periodic
renewal process.

\begin{lemma}[Renewal coordinates]\label{lem:renewal}
Writing $\rho=(1+\mathbb EG)^{-1}$, we have $4\leq\mathbb EG\leq5$
and, for $m\geq0$,
\begin{align}
 \nu_b([0^a1\,0^{g_1}1\cdots0^{g_m}1\,0^b])
   &=\rho\,\Pr(G\geq a)\prod_{j=1}^m p_{g_j}\Pr(G\geq b),
       \label{eq:renewalword}\\
 \nu_b([0^n])&=\rho\sum_{j\geq n}\Pr(G\geq j). \label{eq:zeros}
\end{align}
In particular the denominator in
\begin{equation}\label{eq:gapratio}
 R_k(\nu)=
 \frac{\nu([1\,0^{2k}1])}
 {\nu([1\,0^{2k}1])+\nu([1\,0^{2k+1}1])}
\end{equation}
is $\rho2^{-k}\geq2^{-k}/6$ at $\nu_b$, and
$R_k(\nu_b)=b_k$. The measure $\nu_b$ is limit-computable,
uniformly from the approximation program.
\end{lemma}
\begin{proof}
Each renewal block has one one and $G$ zeros. Its mean length is
$1+\mathbb EG$, so renewal starts have intensity $\rho$. Conditional
on a start, forward and backward gaps have their original independent
laws. This proves \eqref{eq:renewalword}. There are
$(G-n+1)_+$ possible starts of an all-zero word of length $n$ within
a gap, giving \eqref{eq:zeros}. Also
\[
 \mathbb EG=\sum_{k\geq1}2^{-k}(2k+1-b_k)\in[4,5].
\]
The even and odd specified gaps have total mass $2^{-k}$,
which proves the ratio assertions, including $b_k=0,1$.

For effective tails use
\begin{equation}\label{eq:renewaltail}
 \sum_{k>K}2^{-k}(2k+1)=(2K+5)2^{-K}.
\end{equation}
Thus every cylinder coordinate is uniformly $0'$-computable.
More explicitly, approximate each rational code $b_k$ by a total
computable array of rational guesses, clamped to $[0,1]$, eventually
correct for each fixed $k$. At stage $s$ use the first $s$ gap
pairs and normalize their total $1-2^{-s}$. Their stationary
finite-state renewal measures are computable. For every fixed
finite prefix of pairs the guesses settle, while
\eqref{eq:renewaltail} uniformly controls the remaining mean and
cylinder errors. These computable measures converge to $\nu_b$.
The finite algorithm contains no oracle for $0'$.
\end{proof}

Mixing is not needed and need not hold. If every $b_k=0$, all
interarrival lengths $G+1$ are even. The intensity is $1/6$, but
$\nu_b(x_0=x_{2n+1}=1)=0$ at every odd separation. The relevant
period is that of interarrival lengths, not zero-run lengths.

\subsection{The two realization inputs}
Hellouin de Menibus--Sablik \cite[Corollary~1]{hdms} prove that
every limit-computable shift-invariant measure on a finite output
alphabet is the limit of a CA $F$ on an enlarged alphabet $A$,
from every full-support shift-ergodic measure on $A^\Z$.
The target is supported on the output subalphabet; auxiliary symbols
vanish in the limit. This is the singleton target-measure result;
no stronger seed-mixing hypothesis for general limit sets is used.
This seed quantifier concerns the pre-blocking alphabet $A$
and its one-cell shift. Write $\mu_{\rm row}$ for such an $A$-row
law. The physical preparation clones its blocked image on
$\Sigma=A^B$. The canonical uniform Bernoulli row on $\Sigma$
arises by blocking the uniform Bernoulli row on $A$.
No claim is made here for arbitrary ergodic physical block-row laws.
The self-organizer formats segments, computes approximating words,
copies them and merges segments that are too short. We import this
finite-rule construction rather than reprove its segment estimates.

Pad the radius of $F$ to an integer $r\geq1$, let $B=2r$, and
write $(Sx)_i=x_{i+1}$. For $H=S^rF$,
\begin{equation}\label{eq:blockdrift}
 H^t(x)_i=F^t(x)_{i+rt}.
\end{equation}
Grouping $B$ original cells turns $H$ into a two-block rule $h$
with neighbourhood $\{0,1\}$ on $\Sigma=A^B$. Before applying
$h$, take full-symbol NEC majority in every redundant plane, with
a fixed tie rule for three distinct symbols. The stacked rule
$G_\beta$ has six parents
\[
 \{(0,0,0),(1,0,0),(0,1,0),
                 (0,0,1),(1,0,1),(0,1,1)\}.
\]
It is an ordinary finite-alphabet TI rule; the finite target
approximation program is part of $F$.

The arbitrary-trajectory reliability theorem of G\'acs--Reif
\cite{gacsreif}, in the explicit form of Taati--Mi\k{e}kisz
\cite[Theorem~3.1]{miekisztaati}, applies to every cloned initial
row and all logical times. For quantitative choices we also use
G\'acs \cite[Theorem~1]{gacstoom}, whose infinite-plane site-error
bound is $4D_*\varepsilon$ for
$D_*=24$ and $\varepsilon<1/(32D_*^8)$. The factor $D_*$
is part of that bound. These results descend from Toom's
stability construction \cite{toom}.

Fix the rational number
\begin{equation}\label{eq:epsilon}
 \varepsilon_*=\frac1{64D_*^8},\qquad
 U_\beta=\{K:K(a\mid r)>0,\quad
          K(G_\beta(r)\mid r)>1-\varepsilon_*\ \text{for all }a,r\}.
\end{equation}
Uniform positive rational wrong-symbol noise of total mass
$\varepsilon_*/2$ gives a computable member. The set is nonempty
and relatively open. Conditional probabilities of actual faults
are at most $\varepsilon_*$; conditioning backwards through time
layers gives the finite joint-fault bound required by reliability.
Actual fault indicators need not be unconditionally independent.
The one-site error is at most $96\varepsilon_*<1/20$.

\subsection{Age-dependent amplification}
At integer $t\geq2$ use the odd integer schedule
\[
 m_t=2\lceil\log_2t\rceil^2+1
\]
and copies spaced by $2t+1$ in the $a$ direction. This integer
rounded log-squared schedule is computed by bit length and has
the required growth and summability. Decode a blocked symbol by
strict majority, using a fixed default if no symbol has a majority,
and then unpack the $B$ original coordinates. These operations are
one schema on labelled block alphabets.

Conditional on the whole seed row $Y=y$, the initial data are
deterministic. Separated endpoint cones use disjoint update coins,
so their symbols are conditionally independent. Each has the same
intended symbol and error at most $1/20\leq1/4$. Hoeffding gives
\begin{equation}\label{eq:age-error}
 \Pr(\text{wrong decoded block}\mid Y=y)\leq e^{-m_t/8}=:e_t.
\end{equation}
For a fixed original word $u$, let $W$ be the finitely many
blocks covering it. A union bound and integration over $y$ yield
\begin{equation}\label{eq:worderror}
 \left|\Pr_{\lambda P_K^t}(\text{decoded }u)
                      -(F_*^t\mu_{\rm row})([u])\right|\leq |W|e_t.
\end{equation}
Here $\mu_{\rm row}$ is the pre-blocking $A$-row law.
Its one-cell shift invariance passes to $F_*^t\mu_{\rm row}$,
so \eqref{eq:blockdrift} causes no moving-frame correction:
that invariance absorbs the drift $rt$ on the original cells.
Block-shift invariance alone would not justify this step.
There is still age in the copy geometry.

\begin{proof}[Proof of the depth-one part of Theorem~\ref{thm:realization}]
Apply \eqref{eq:worderror} to the two gap words in
\eqref{eq:gapratio}. The noiseless probabilities converge to
their renewal values by the self-organization input. Their
limiting sum is $\rho2^{-k}>0$, so the inner time limit of the
decoded ratio exists and equals $b_k$ for every fixed $k$.
Every finite-time denominator is positive: a finite assignment of
the copies can realize either binary word, and positive-rate
updates give it positive probability. The initial full-support
cloned row also gives these events positive probability.
Taking $k\to\infty$ gives $\beta$. The bound is uniform over
$U_\beta$ and over row contents. The noiseless theorem applies to every full-support,
one-cell-shift-invariant ergodic law on $A^\Z$. Thus the same
value holds for clones of their blocked images, including
noncomputable admitted row laws. It is not asserted for every
shift-ergodic law on $\Sigma^\Z$. For the
arithmetic upper bound and later inquiry statement, use the
canonical computable row. Theorem~\ref{thm:ceilings} gives the
reverse inclusion. The image is over finite rules and their
labelled block presentations, not one universal alphabet.
\end{proof}

\subsection{Attainment at every finite depth}
\begin{proof}[Proof of the remaining part of Theorem~\ref{thm:realization}]
For $d\geq1$ and $\beta\leq_T0^{(d+1)}$, iterated relativized
real limit approximation gives rationals
$b_{k_1,\ldots,k_d}\in[0,1]$, uniformly $0'$-computable, whose
$d$ prescribed nested limits exist and equal $\beta$. This uses
real Cauchy names and uniform limit approximation, not uniformly
decidable strict cuts of arbitrary oracle reals. Fix a computable
bijection $p_d$ from positive $d$-tuples to positive integers.
Put $b_{\mathbf k}$ into gap pair $p_d(\mathbf k)$ in
\eqref{eq:gaplaw}. The same geometric bound controls all tails,
so the renewal measure is still limit-computable. Use the same
uniform readout schema, now asking gap question $p_d(\mathbf k)$.
Its inner time limit is $b_{\mathbf k}$ and the $d$ outer limits
give $\beta$. Every intermediate limit is part of the hypothesis
on that chosen array; arbitrary programs are not presumed to
converge.

At $d=0$, choose a limit-computable Bernoulli target with parameter
$\beta\in\D_2\cap[0,1]$. Its cylinder probabilities are limits of
computable probabilities by a computable rational approximation to
$\beta$. Apply the same noiseless realization and reliable stack,
then read the decoded one-site probability with \eqref{eq:age-error}.
Its single time limit is $\beta$. Endpoints zero and one are
allowed target measures even though the initial row has full support.
The ceilings in Theorem~\ref{thm:ceilings} give exactness at each
depth. The schema is fixed separately for each $d$.
\end{proof}
